\section{Front-Loading 的五条实验课}

下面五张图不是重复证明同一件事，而是逐步回答五个反对意见：预训练后立即有效吗？SFT 会洗掉收益吗？高质量数据是否过拟合？更多 SFT compute 或更多后期数据能追上吗？完成 RLVR 后还剩多少？每张图都必须先确认比较对象，再读柱状高度。

\subsection{Lesson 1：预训练后立即出现收益}

第一条 lesson 先检查最直接的时间点：尚未进入 SFT 或 RL 时，提前加入 reasoning data 是否已经改变 base model。这个检查把后训练因素全部排除，只观察预训练阶段的数据分配是否足以形成可测差异。Reason base 与 no-reason base 在完成 pretraining 后比较，Ampere 的平均相对提升约为 16\%。这说明 reasoning data 不只是教模型输出某种格式，而可能改变 base representation 对数学、科学与通用推理任务的可用性。它仍不证明模型内部产生了忠实可解释的 reasoning process。

\lecturefigure{slide-017.jpg}{Lesson 1：在预训练阶段加入 reasoning data 的即时收益}{Stanford 官方录像，00:15:48--00:16:36}

\begin{knowledgebox}{读图：比较顺序}
先确认两组模型在 architecture、token budget 与普通预训练 mixture 上一致；再看每个 benchmark 的蓝绿柱；最后看平均相对提升。若某一任务起点很低，小的绝对变化会形成较大的相对百分比，因此需要同时看 raw score 与 relative gain。
\end{knowledgebox}

\subsection{Lesson 2：SFT 没有把优势洗掉}

有即时收益还不够，因为生产模型通常都会继续接受 SFT；因此第二条 lesson 把相同后训练施加到两种 foundation 上，检查优势是否只是短暂现象。两组 base model 接受相同 SFT 后，reason base 仍保持约 9.3\% 平均相对优势。这里的教学结论是“foundation matters”：后训练不是完全重置模型，而是在 base representation 上继续优化。早期收益可能缩小、保持或在特定任务上放大，必须逐任务检查。

\lecturefigure{slide-018.jpg}{Lesson 2：早期 reasoning 收益在 SFT 后仍然存在}{Stanford 官方录像，00:16:36--00:17:14}

\teachervoice{00:16:36--00:17:14，讲者反复使用“washed away”这个说法。她关心的不只是 base score，而是同一优势在后训练之后是否仍可观测；这是 front-loading 是否值得工程投入的关键判据。}

\subsection{Lesson 3：高质量数据可能具有 latent effect}

第三条 lesson 转向一个更隐蔽的问题：若高质量数据在 base evaluation 上没有立刻抬高分数，是否意味着它没有价值？课程把推理数据分为 SHQ、LDQ 与 LMQ。SHQ 小而高质量，LDQ 大而多样但质量较低，LMQ 为二者组合。预训练结束时，LMQ 与 LDQ 可能几乎相同；经过 SFT 后，LMQ 才显示约 4.25\% 的额外提升。高质量信号并非立即反映在 base benchmark 上，而可能为后续优化提供更好的可塑性。

\lecturefigure{slide-019.jpg}{Lesson 3：高质量预训练数据的潜在收益在 SFT 后被激活}{Stanford 官方录像，00:17:14--00:18:48}

\begin{knowledgebox}{术语表：SHQ、LDQ、LMQ}
\begin{tabularx}{\textwidth}{L{0.16\textwidth}L{0.23\textwidth}X}
\toprule
缩写 & 数据属性 & 教学作用 \\
\midrule
SHQ & Small, High Quality & 检验少量精筛数据是否提供额外结构 \\
LDQ & Large, Diverse, lower Quality & 检验规模与多样性基础 \\
LMQ & SHQ + LDQ 的组合 & 检验高质量成分是否在后训练后显现 latent effect \\
\bottomrule
\end{tabularx}
\end{knowledgebox}

\subsection{Lesson 4：后期加倍 compute 或挪动数据仍追不上}

左图给 no-reason base 更多 SFT epochs，右图固定 reasoning data 总预算并比较“全放 SFT”与“pretraining + SFT 分配”。两种情况下，front-loaded model 仍保持优势。这分别反驳“只是后训练不够久”和“只要把同样数据晚点用就行”两个解释。

\lecturefigure{slide-020.jpg}{Lesson 4：compute-matched 与 data-matched 条件下的强地基效应}{Stanford 官方录像，00:18:48--00:20:54}

\begin{importantbox}{为什么要同时做两种匹配}
Compute matching 控制训练计算量，data matching 控制样本预算；二者解决不同混杂因素。若 early condition 既看了更多数据又花了更多 FLOPs，就不能把收益归因于时机。严谨报告至少应列 token、epochs、optimizer steps、sequence length 与 rollout/SFT 的额外成本。
\end{importantbox}

\subsection{Lesson 5：经过 RLVR 后仍有 durable advantage}

完成 SFT 后仍领先，依然不能保证 RLVR 不会把两条路径重新拉平；第五条 lesson 因此检查最终可交付 pipeline，而不是停在中间 checkpoint。Ampere 相对 Volta 的平均相对优势约 19\%，AIME 等高难数学任务的相对差异更大。图支持“early reasoning 可以产生 durable and compounding advantage”，但不说明所有任务都会同幅增长，也不说明后训练不再重要；恰恰是相同后训练把不同 foundation 的差异暴露出来。

\lecturefigure{slide-021.jpg}{Lesson 5：Front-loading 在完整 post-training 后仍保持优势}{Stanford 官方录像，00:20:54--00:21:48}

\begin{warningbox}{大相对提升不等于接近满分}
AIME 等任务基线较低时，几分的绝对提高可能对应很大的 relative gain。工程决策应同时报告绝对分数、平均口径、试验次数和方差；不能只截取最大的相对百分比作为通用能力提升。
\end{warningbox}

\subsection{课堂问答补全了哪些论文外的工程细节}

中场问答说明，reasoning data 使用社区已有数据集与工作定义，常见样本是 question、long reasoning trace、final answer，且实验组成偏 STEM。Quality 可以由过滤强度、复杂度、长度或 difficulty label 操作化；recipe optimization 是自动估计与人工组装的混合流程。讲者还明确说，反转 phase 顺序在其消融中更差。

\teachervoice{00:22:00--00:27:58，讲者没有把“reasoning”包装成跨领域统一概念。她承认所用数据偏数学、代码与 STEM，最终 blend 也并非全自动。这些限制决定了结果的外推边界。}

\begin{knowledgebox}{实践经验：如何记录一个可复现 data recipe}
至少保存数据源版本、去重键、过滤器与阈值、quality rubric、domain mapping、unique token 数、目标 epochs、每阶段权重、phase boundary、随机种子与 proxy-evaluation 配置。只发布最终百分比，无法区分质量估计、采样实现或顺序造成的差异。
\end{knowledgebox}

\subsection{本章小结}

五条 lesson 建立了 front-loading 的证据链：收益立即出现、SFT 后不消失、高质量信号可能延迟显现、更多后期 compute/data 不能完全追平、完整 RLVR 后优势仍在。下一部分进一步追问：除了更早提供 reasoning data，能否让模型在预训练时主动探索自己的 thought？

